\documentclass[11pt]{article}
\usepackage[margin=1in]{geometry}
\usepackage{amsmath,amssymb,amsthm,mathtools}
\usepackage{graphicx}
\usepackage{xurl}
\usepackage[colorlinks=true,linkcolor=blue,citecolor=blue,urlcolor=blue]{hyperref}
\newtheorem{theorem}{Theorem}[section]
\newtheorem{lemma}[theorem]{Lemma}
\newtheorem{corollary}[theorem]{Corollary}
\newtheorem{remark}[theorem]{Remark}
\newcommand{\Tr}{\operatorname{Tr}}
\newcommand{\Aav}{\mathcal A}
\newcommand{\EA}{E_A}
\newcommand{\C}{\mathbb C}
\emergencystretch=1em
\makeatletter
\renewcommand{\@seccntformat}[1]{\csname the#1\endcsname\hspace{0.8em}}
\makeatother
\title{Low-Dimensional, Rank-One, and First Higher-Cycle\\
Closures of the Pinching Refinement Conjecture}
\author{Automated mathematics run \texttt{fa\_banach\_001}, lane 3}
\date{11 August 2026}

\begin{document}
\maketitle
\begin{abstract}
Dinh asked whether the normalized trace average of all words with \(n\)
letters \(A\) and \(m\) letters \(B\) dominates the same average after
pinching \(B\) onto the commutant of \(A\). We prove the conjecture
(i) for all positive semidefinite \(2\times2\) matrices and all
\(n,m\geq0\), (ii) for rank-one \(A\) in every dimension and all
\(n,m\geq0\), (iii) in every dimension when \(n\leq1\), and
(iv) in every dimension at the first higher-cycle corner
\((n,m)=(2,3)\). A more general path argument covers every pair for which,
in an eigenbasis of
\(A\), a diagonal unitary gauges \(B\) to an entrywise nonnegative
matrix. The general higher-dimensional, cyclic-phase case remains open.
The status is \emph{candidate partial result, likely valid, pending human review}.
\end{abstract}

\section{Source question and notation}
The source is Trung Hoa Dinh, \emph{On the Failure of the Upper Bound in the
Refined BMV Conjecture and a Pinching Correction}, arXiv:2605.17782 (2026)
\cite{Dinh}. Conjecture~5.1, on PDF page~7, is reproduced in
Figure~\ref{fig:source}. The source proves the conjecture when \(m=2\), but
leaves the general case open.
\begin{figure}[ht]
  \centering
  \includegraphics[width=0.95\textwidth]{figures/open_problem_crop.png}
  \caption{Conjecture~5.1 and its equivalent form in the source paper,
  cropped from PDF page~7.}
  \label{fig:source}
\end{figure}

For positive semidefinite matrices \(A,B\in M_d(\C)\), let
\begin{equation}\label{eq:average}
 \Aav_{n,m}(A,B)
 =\frac{1}{\binom{n+m}{n}}\sum_W \Tr W(A,B),
\end{equation}
where the sum runs over all words containing exactly \(n\) letters \(A\)
and \(m\) letters \(B\). If
\(A=\sum_{\lambda}\lambda P_{\lambda}\) is the spectral decomposition, set
\begin{equation}\label{eq:pinching}
 \EA(B)=\sum_{\lambda}P_{\lambda}BP_{\lambda}.
\end{equation}
Because \(A\) and \(\EA(B)\) commute, the source conjecture is
\begin{equation}\label{eq:conjecture}
 \Aav_{n,m}(A,B)\geq
 \Aav_{n,m}(A,\EA(B))
 =\Tr\!\left(A^n\EA(B)^m\right)
 \qquad(n,m\geq0).
\end{equation}
Zero powers are identities. Thus the claim is: \emph{for all positive
semidefinite \(A,B\) and all nonnegative integers \(n,m\), inequality
\eqref{eq:conjecture} holds.}

\section{Proof intuition}
In an eigenbasis of \(A\), expand the trace of a word as a sum over closed
index paths. An \(A\)-step cannot change the current spectral vertex, while
a \(B\)-step may. Pinching deletes precisely the \(B\)-edges that cross
between distinct eigenspaces of \(A\). Thus the desired difference is the
sum of weights of closed paths that use at least one deleted edge.

This is not a proof in general: in dimension at least three, a closed
\(B\)-cycle can carry a nontrivial complex phase and hence have negative
real part. It does prove the claim when all \(B\)-entries can be made
nonnegative by a diagonal phase gauge. Every \(2\times2\) pair has this
property, because a single off-diagonal phase can always be removed.

Two further arguments avoid path phases. If \(A\) has rank one, cyclicity
rewrites the word average as an average over weak compositions of \(m\);
rank-one collapse then reduces each term to scalar moments of \(B\), which
dominate the corresponding powers of the first moment. If \(n\leq1\), the
same scalar moment inequality applied to compressions of \(B\) proves the
claim in every dimension. At \((n,m)=(2,3)\), cyclicity leaves just two
trace expressions. One is another compression moment; the other becomes a
squared-norm estimate after normalizing \(B\) to a correlation matrix.
None of these mechanisms controls arbitrary higher-dimensional cyclic
phases, so the original full conjecture remains open.

\section{A compression moment lemma}
\begin{lemma}[compression moments]\label{lem:compression}
Let \(B\geq0\), let \(P\) be an orthogonal projection, and let \(r\) be a
nonnegative integer. Then
\begin{equation}\label{eq:compression}
 \Tr(PB^rP)\geq \Tr\bigl((PBP|_{P\C^d})^r\bigr).
\end{equation}
Equivalently, with both sides viewed on \(\C^d\),
\(\Tr(PB^rP)\geq\Tr((PBP)^r)\) for \(r\geq1\); for \(r=0\), the right-hand
side is understood as the identity on \(P\C^d\).
\end{lemma}
\begin{proof}
Diagonalize the positive compression \(C=PBP|_{P\C^d}\), say
\(Ce_j=\mu_j e_j\) for an orthonormal basis of \(P\C^d\). If
\(B=\sum_k t_k Q_k\) is a spectral decomposition, then the nonnegative
numbers \(\langle e_j,Q_ke_j\rangle\) sum to one. Convexity of
\(t\mapsto t^r\) on \([0,\infty)\) gives
\[
 \langle e_j,B^r e_j\rangle
 =\sum_k t_k^r\langle e_j,Q_ke_j\rangle
 \geq
 \left(\sum_k t_k\langle e_j,Q_ke_j\rangle\right)^r
 =\mu_j^r.
\]
(The cases \(r=0,1\) are equalities.) Summing over \(j\) proves the lemma.
\end{proof}

\section{The path-positive class and the \(2\times2\) theorem}
\begin{theorem}[phase-balanced pairs]\label{thm:phase}
Suppose there is an orthonormal eigenbasis of \(A\) and a diagonal unitary
\(D\) in that basis such that \(D^*BD\) is entrywise real and nonnegative.
Then \eqref{eq:conjecture} holds for every \(n,m\geq0\).
\end{theorem}
\begin{proof}
Conjugating both matrices by \(D\) changes neither side of
\eqref{eq:conjecture}; moreover \(D\) commutes with \(A\), and pinching is
covariant under this conjugation. We may assume that \(A\) is diagonal with
nonnegative diagonal entries and every entry of \(B\) is nonnegative.

Fix \(W=X_1\cdots X_{n+m}\), with each \(X_s\) equal to \(A\) or \(B\).
Its trace has the closed-path expansion
\begin{equation}\label{eq:path}
 \Tr W=\sum_{i_0,\ldots,i_{n+m-1}}
 (X_1)_{i_0i_1}(X_2)_{i_1i_2}\cdots
 (X_{n+m})_{i_{n+m-1}i_0}.
\end{equation}
Every summand is nonnegative. If \(X_s=A\), its factor vanishes unless the
path stays at the same vertex. Replacing \(B\) by \(\EA(B)\) sets
\(B_{ij}\) to zero exactly when the eigenvalues of \(A\) at \(i\) and \(j\)
are different, and leaves all other entries unchanged. Hence pinching
deletes a subset of nonnegative summands in \eqref{eq:path}. Thus
\(\Tr W(A,B)\geq\Tr W(A,\EA(B))\) word by word. Averaging proves the claim.
\end{proof}

\begin{corollary}[all \(2\times2\) cases]\label{cor:two}
For all positive semidefinite \(A,B\in M_2(\C)\) and all \(n,m\geq0\),
inequality \eqref{eq:conjecture} holds.
\end{corollary}
\begin{proof}
Choose an eigenbasis of \(A\). If \(A\) is scalar, then \(\EA(B)=B\).
Otherwise write
\[
 B=\begin{pmatrix}x&z\\ \overline z&y\end{pmatrix},
 \qquad x,y\geq0.
\]
A diagonal unitary makes \(z\) real and nonnegative, so
Theorem~\ref{thm:phase} applies. Directly, a closed two-vertex path uses
equally many \(1\to2\) and \(2\to1\) transitions, whose product is
\(|z|^{2q}\geq0\).
\end{proof}
\begin{remark}
Repeated eigenvalues cause no problem: pinching retains all edges within
each eigenspace and deletes only cross-eigenspace edges.
\end{remark}

\section{Rank-one \(A\) in arbitrary dimension}
\begin{lemma}[cyclic composition formula]\label{lem:composition}
For \(n\geq1\),
\begin{equation}\label{eq:composition}
 \Aav_{n,m}(A,B)
 =\frac{1}{\binom{n+m-1}{m}}
 \sum_{r_1+\cdots+r_n=m}
 \Tr\bigl(AB^{r_1}AB^{r_2}\cdots AB^{r_n}\bigr),
\end{equation}
where the sum is over weak compositions \(r_j\geq0\).
\end{lemma}
\begin{proof}
Words beginning with \(A\) are in bijection with the weak compositions in
\eqref{eq:composition}. Mark one of the \(n\) occurrences of \(A\) in each
unrestricted word and cyclically rotate that occurrence to the first
position. If \(T\) is the sum over all words and \(S\) the sum over words
beginning with \(A\), this marked rotation is \(N\)-to-one, where
\(N=n+m\), even for periodic words (the marked position is retained).
Cyclic invariance therefore gives \(nT=NS\). Since
\(\frac nN\binom Nn=\binom{N-1}{n-1}=\binom{n+m-1}{m}\), the unrestricted
average equals the average over words beginning with \(A\), proving the
formula.
\end{proof}

\begin{theorem}[rank-one closure]\label{thm:rankone}
If \(A\geq0\) has rank one, then \eqref{eq:conjecture} holds in every finite
dimension for all \(n,m\geq0\).
\end{theorem}
\begin{proof}
The case \(n=0\) follows from Lemma~\ref{lem:compression}, summed over the
spectral projections of \(A\), and is also in Theorem~\ref{thm:smalln}.
Assume \(n\geq1\), and write \(A=aP\), where \(a>0\) and
\(P=|v\rangle\langle v|\) for a unit vector \(v\). Put
\(\mu=\langle v,Bv\rangle\). Rank-one collapse in \eqref{eq:composition}
gives
\begin{equation}\label{eq:rankcollapse}
 \Tr\bigl(AB^{r_1}\cdots AB^{r_n}\bigr)
 =a^n\prod_{j=1}^n\langle v,B^{r_j}v\rangle.
\end{equation}
For every nonnegative integer \(r\), the scalar spectral-measure argument in
Lemma~\ref{lem:compression} gives
\(\langle v,B^rv\rangle\geq\mu^r\). Therefore every summand in
\eqref{eq:composition} is at least \(a^n\mu^m\), so
\begin{equation}\label{eq:rank-lhs}
 \Aav_{n,m}(A,B)\geq a^n\mu^m.
\end{equation}
The spectral projections of \(A\) are \(P\) and \(I-P\), hence
\(\EA(B)=\mu P+(I-P)B(I-P)\), and consequently
\[
 \Tr\bigl(A^n\EA(B)^m\bigr)=a^n\mu^m.
\]
Together with \eqref{eq:rank-lhs}, this proves the claim.
\end{proof}
\begin{remark}
Rank-one \(A\) also follows from Theorem~\ref{thm:phase}: diagonalize the
compression of \(B\) to \(v^\perp\), then adjust eigenvector phases so that
all entries coupling them to \(v\) are nonnegative. The composition proof
gives an independent scalar-moment mechanism.
\end{remark}

\section{The range \(n\leq1\) in arbitrary dimension}
\begin{theorem}[zero or one \(A\)-letter]\label{thm:smalln}
For arbitrary positive semidefinite matrices \(A,B\) in every finite
dimension, \eqref{eq:conjecture} holds for \(n=0\) and \(n=1\), for every
\(m\geq0\).
\end{theorem}
\begin{proof}
Let \(A=\sum_\lambda\lambda P_\lambda\), and write
\(C_\lambda=P_\lambda BP_\lambda|_{P_\lambda\C^d}\).
If \(n=0\), Lemma~\ref{lem:compression} gives
\[
 \Aav_{0,m}(A,B)=\Tr(B^m)
 =\sum_\lambda\Tr(P_\lambda B^mP_\lambda)
 \geq\sum_\lambda\Tr(C_\lambda^m)
 =\Tr(\EA(B)^m).
\]
If \(n=1\), cyclic invariance makes every word trace equal to
\(\Tr(AB^m)\). Applying Lemma~\ref{lem:compression} blockwise and weighting
by the nonnegative eigenvalues \(\lambda\) yields
\begin{align*}
 \Aav_{1,m}(A,B)
 &=\Tr(AB^m)
 =\sum_\lambda\lambda\Tr(P_\lambda B^mP_\lambda)\\
 &\geq\sum_\lambda\lambda\Tr(C_\lambda^m)
 =\Tr\bigl(A\EA(B)^m\bigr).
\end{align*}
\end{proof}

\section{The first higher-cycle corner}
\begin{theorem}[the case \((n,m)=(2,3)\)]\label{thm:twothree}
For arbitrary positive semidefinite matrices \(A,B\) in every finite
dimension,
\[
 \Aav_{2,3}(A,B)\geq \Tr\bigl(A^2\EA(B)^3\bigr).
\]
\end{theorem}
\begin{proof}
Choose an eigenbasis of \(A\), and within each eigenspace choose a basis
that diagonalizes the corresponding compression of \(B\). Then
\(D=\EA(B)\) is diagonal in this basis. The cyclic-composition formula,
now marking one of the three \(B\)-letters, gives
\[
 \Aav_{2,3}(A,B)
 =\frac{1}{6}\sum_{p+q+r=2}
   \Tr\bigl(BA^pBA^qBA^r\bigr)
 =\frac12\left\{\Tr(A^2B^3)+\Tr(ABAB^2)\right\}.
\]
Lemma~\ref{lem:compression}, applied on every spectral block of \(A\) and
weighted by the corresponding eigenvalue squared, yields
\begin{equation}\label{eq:first-23}
 \Tr(A^2B^3)\geq\Tr(A^2D^3).
\end{equation}

It remains to prove the same lower bound for the second trace. Write
\(d_i=b_{ii}\). If \(d_i=0\), positivity of \(B\) forces its \(i\)-th row
and column to vanish, so those indices may be discarded. On the remaining
indices put
\[
 C=D^{-1/2}BD^{-1/2},
 \qquad X=AD.
\]
Then \(C\geq0\), \(c_{ii}=1\), and \(X=\operatorname{diag}(x_i)\geq0\)
with \(x_i=a_i d_i\). Direct substitution and cyclicity give
\begin{align}
 \Tr(ABAB^2)
 &=\Tr(XCXC DC)\notag\\
 &=\sum_k d_k\,(CXCXC)_{kk}.\label{eq:corr-reduction}
\end{align}
For each \(k\),
\[
 (CXCXC)_{kk}
 =\left\|C^{1/2}XC e_k\right\|^2.
\]
Since \(\|C^{1/2}e_k\|^2=c_{kk}=1\), Cauchy--Schwarz gives
\begin{align*}
 \left\|C^{1/2}XC e_k\right\|^2
 &\geq
 \left|\left\langle C^{1/2}e_k,C^{1/2}XC e_k\right\rangle\right|^2\\
 &=\left(\sum_j x_j|c_{jk}|^2\right)^2
 \geq x_k^2.
\end{align*}
Substitution in \eqref{eq:corr-reduction} proves
\[
 \Tr(ABAB^2)\geq\sum_k d_kx_k^2
 =\sum_k a_k^2d_k^3
 =\Tr(A^2D^3).
\]
Combining this with \eqref{eq:first-23} and the displayed average identity
proves the theorem.
\end{proof}

\section{Verification and adversarial checks}
The script \texttt{code/check\_pinching\_subcases.py} independently computes
word averages as coefficients of \(\Tr(A+tB)^{n+m}\), constructs spectral
pinching, and checks random complex instances of all proved classes. It is a
sanity check, not part of any proof. It covers random \(2\times2\) pairs,
arbitrary-dimensional \(n=0,1\), random rank-one \(A\), phase-balanced
pairs, the all-dimensional \((n,m)=(2,3)\) theorem, and the source's known
\(m=2\) range.

A bounded counterexample search preceded promotion: random real and complex
PSD matrices in dimensions \(3,4,5\), ranks from one to full, total word
degree up to \(20\), continuous optimization over Gram factors, and
structured correlation-matrix families. No negative gap was found. This is
evidence only. Convexity in \(B\) was tested and rejected as a proof route:
the Hessian of the full coefficient can have negative directions, so a
direct Jensen argument for the full word average is unavailable.

The formal proofs have these independent failure checks.
\begin{itemize}
\item Theorem~\ref{thm:phase} keeps, rather than deletes, edges inside
repeated eigenspaces.
\item In Corollary~\ref{cor:two}, Hermitian symmetry makes the two directed
off-diagonal factors conjugates.
\item The normalization in \eqref{eq:composition} is exactly the number
\(\binom{n+m-1}{m}\) of weak compositions.
\item Lemma~\ref{lem:compression} uses scalar convexity in spectral measures,
not the false claim that \(t^r\) is operator convex for every \(r\).
\item Theorem~\ref{thm:twothree} removes zero diagonal entries before using
\(D^{-1/2}\), and its correlation estimate is a squared norm rather than a
termwise assertion about signed three-cycles.
\end{itemize}

\section{Scope, novelty, and review}
This packet does \emph{not} prove Conjecture~5.1 in full. It closes the first
higher-cycle corner \((n,m)=(2,3)\), but leaves open the remaining cases
with \(d\geq3\), \(n\geq2\), \(m\geq3\) and genuinely unbalanced cycle
phases. The source already proves \(m=2\) in all dimensions; that known case
is not claimed as new.

A bounded novelty search on 11 August 2026 used the exact arXiv id, title,
the phrases \emph{pinching refinement} and \emph{BMV pinching conjecture},
and exact text from Conjecture~5.1. It checked the local run indexes and
arXiv/web results. It found the source and literature on the failed
refined-BMV upper bound, but no later proof, counterexample, or statement of
the \(2\times2\), rank-one-\(A\), \(n\leq1\), phase-balanced, or
\((n,m)=(2,3)\) closures.
This bounded search is not a guarantee of global novelty.

Recommended human review: \emph{promote as a likely-valid partial result
after checking the path-deletion identification in Theorem~\ref{thm:phase}
the marked-word count in Lemma~\ref{lem:composition}, and the correlation
normalization in Theorem~\ref{thm:twothree}.} The genuinely open next target
is control of signed higher-dimensional cycle phases.

\begin{thebibliography}{9}
\bibitem{Dinh}
T.~H. Dinh,
\emph{On the Failure of the Upper Bound in the Refined BMV Conjecture and a
Pinching Correction}, arXiv:2605.17782 (2026),
\url{https://arxiv.org/abs/2605.17782}.
\end{thebibliography}
\end{document}
